\chapter{Architecture}
\label{ch:arch}

As shown in Section \ref{ch:afl}, traditional edge coverage relies on a single shared memory bitmap and a hashing function to track execution paths. This approach is highly performant and has led to thousands of bugs being found in recent years \cite{ossfuzz}. 

However, indirect jumps are generally not explicitly tracked and still present a challenge. While these dynamic jumps can be instrumented at compile time, they inherently represent a 1-to-$N$ control-flow mapping, where a single call site may resolve to numerous distinct targets at runtime. 

Attempting to capture this execution within a traditional, fixed-size coverage map significantly increases the probability of collisions. Alternatively, in collision-free modes, it would inflate the size of the map. Either scenario potentially degrades the fuzzer's efficacy. 

In this chapter, we describe our solution, which addresses these problems by introducing a secondary coverage map dedicated exclusively to tracking indirect jumps.

\section{The Indirect Branch Problem}

As detailed in Section \ref{sub:shm}, AFL++ relies on a highly optimised, fixed-size shared memory bitmap (typically 64KB, although resizeable) to track execution paths. 

Indirect branches --- such as function pointers, virtual method dispatches in C++, and complex switch-case statements --- introduce a dynamic $1$-to-$N$ mapping into the CFG. A single indirect call site $A$ does not resolve to a single destination $B$. Instead, depending on the program state, $A$ can transition to one of many targets during execution.

If these indirect branches are instrumented using the standard edge-tracking mechanism, every new target $B_i$ generates a new edge $A \rightarrow B_i$.

\paragraph{The Collision Problem} Mapping these high-volume dynamic edges into the same fixed-size map used for direct edges can significantly increase the overall density of the coverage map. A denser map naturally increases the probability of collisions. 

When a collision occurs, it causes "coverage blindness". The fuzzer might observe an execution that reaches a novel path, but if that edge hashes to a byte already saturated by an indirect branch, the fuzzer will incorrectly assume the path is redundant. Consequently, viable inputs are discarded.

\paragraph{Collision-Free Instrumentation Issues} This problem persists even when utilizing collision-free instrumentation modes, such as LLVM PCGUARD (Section \ref{sub:llvmmode}). These modes achieve zero collisions by assigning a unique ID to every basic block and edge discovered at compile time, sizing the coverage map exactly to the number of known edges.

Because indirect branches resolve at runtime, they circumvent the static counting. If unpredictable dynamic edges were naively introduced into a collision-free map, they would risk breaking the zero-collision guarantee.

Simply extending the coverage map to accommodate all dynamic edges would risk allocating an excessively large map, introducing performance bottlenecks and causing CPU cache thrashing.

\subsection{Requirements for a Solution}

To accurately track indirect branches without polluting the primary coverage map or degrading execution speed, an alternative data structure is required. This structure must satisfy two conflicting constraints: 
\begin{itemize}
    \item \textit{High capacity for dynamic edges:} It must be able to absorb a high volume of unpredictable, runtime-generated edges.

    \item \textit{Strict memory and performance bounds:} It must maintain a small memory footprint to avoid cache thrashing and preserve fuzzing throughput.
\end{itemize}

These requirements point towards a probabilistic data structure capable of tracking set membership with a fixed size and a predictable collision rate --- a role perfectly suited for Bloom filters.

\section{The Secondary Indirect Coverage Map}

To accurately track indirect branches without polluting the primary coverage map or degrading execution speed, an alternative data structure is required. We propose the introduction of a dedicated secondary coverage map for tracking indirect control flow transfers. This architecture aims to resolve the conflicting constraints of maintaining high capacity for dynamic edges while preserving strict memory and performance bounds.

The primary concept relies on utilizing a small Bloom filter for each possible source of an indirect jump: each target reached is then represented by a single bit raised to 1 within the corresponding filter. 

\paragraph{Isolation} We aim to decouple the tracking of indirect branches from the standard coverage map and move this data into an independent shared memory region. Isolation ensures that collisions occurring within the secondary map cannot inadvertently suppress standard edge discovery. 

\paragraph{Collision-Free Identification} For our secondary map, we use a dynamic, sequential "\textit{Guard Array}" architecture. In this model, every indirect branch source in the compiled binary is assigned a guaranteed, sequentially numbered and collision-free slot. 

During the compilation phase, a dedicated ELF section is created, allocating exactly one slot for every \texttt{call} or \texttt{jmp} instruction that resolves dynamically. This approach guarantees that the coverage map possesses a 1-to-1 mapping for the \textit{origins} of all indirect control flow transfers.

\paragraph{Configurable Fidelity} To address the diverse requirements of different target applications, we provide configurable fidelity. The size of the Bloom filter, and consequently the collision tolerance of the secondary map itself, can be tuned. This allows balancing the memory footprint against the granularity of fuzzing feedback.

\section{Compile-Time Instrumentation}

As already detailed (Section \ref{sec:compandinst}), visibility into the target application is achieved through compile-time instrumentation. In our solution, the instrumentation logic is integrated into the LLVM compiler infrastructure, specifically extending the existing LLVM mode's PCGUARD pass. 

The instrumentation process is explicitly designed for indirect branches and consists of three primary phases: target identification, metadata allocation, and callback injection.

\paragraph{Target Identification} The first phase of the instrumentation pass involves traversing the LLVM Intermediate Representation to identify instructions capable of dynamic control flow transfers. 

During the traversal of each basic block, the pass filters the instruction stream to isolate two specific LLVM IR instructions: 
\begin{itemize}
    \item \texttt{IndirectBrInst}: Represents an indirect branch instruction --- e.g., computed \texttt{goto} statements and jump tables.

    \item \texttt{CallInst}: Represents a function call. In this case, the pass further verifies that the target operand is an indirect function pointer call or a virtual method dispatch.
\end{itemize}

To maintain feature parity with the broader AFL++ ecosystem, this identification phase respects the coverage boundaries set for standard edge coverage. Fuzzing campaigns can utilise allowlists and denylists to constrain instrumentation to specific source files or functions.

\paragraph{Guard Array Architecture} Once an indirect instruction has been verified and selected for instrumentation, the compiler allocates a slot to track its execution. To track each indirect branch source, our architecture utilises a metadata structure referred to as "Guard Array".

For every identified indirect instruction, the compiler pass allocates a 32-bit integer slot --- a \texttt{GuardPtr} --- into a dedicated contiguous ELF section named "\texttt{\_\_sancov\_indir\_guard}" (mimicking the operation of PCGUARD mode).

As the compiler processes each translation unit, it continuously appends 32-bit slots into this section. The resulting linked binary will contain a sequential array of uninitialized metadata slots, corresponding one-to-one to every instrumented indirect branch source. At runtime, these slots will be populated with unique sequential identifiers.

\paragraph{Callback Injection} The final phase of static instrumentation involves modifying the target's executable code to report the control flow transition. This requires injecting a C function callback into the LLVM IR.

Modifying the IR of a basic block while actively iterating over its instructions risks invalidating iterators and triggering undefined behaviour within subsequent analysis passes. To circumvent this, a two-phase approach is taken: 
\begin{enumerate}
    \item \textit{Collection Phase:} The pass first traverses the CFG, applying the identification and allowlist logic to build a set of valid targets.

    \item \textit{Mutation Phase:} The pass performs the necessary IR modification on the collected set of targets.
\end{enumerate}

For each target instruction in the mutation set, the compiler injects a call to \texttt{\_\_afl\_trace\_indir} immediately preceding the indirect branch. The callback takes two arguments: the address of the corresponding \texttt{GuardPtr} slot and the target address of the indirect branch, which the program has just evaluated.

\section{Runtime Support}

The runtime component of the indirect coverage system operates alongside the standard compiler runtime and serves as a bridge between the metadata injected by the compiler pass and the actual indirect coverage tracking mechanism of the fuzzer.

Its responsibilities are manifold, spanning from the initialisation of the indirect guard array to the execution of the coverage callback.

\paragraph{Guard Array Initialisation} The lifecycle of the indirect tracking mechanism begins before the target program's \texttt{main} function is executed. The LLVM compiler pass registers a module-level constructor that invokes the initialisation routine.

The objective of this initialisation phase is to transition the guard slots from their uninitialized state into distinct identifiers. The function implements a linear scan across the region of the Guard Array dedicated to the corresponding module. A globally-scoped counter is increased monotonically for each indirect jump source within each module.

If Dynamic Shared Objects (DSOs) are loaded into the process space --- and are instrumented --- their constructors invoke the initialisation function within the bounds of their specific \texttt{\_\_sancov\_indir\_guards} section. The runtime incorporates these new slots into the identifier namespace, ensuring a contiguous ID space across the main executable and all loaded libraries.

Upon completion of the initialisation scan, the final ID value defines the maximum required index for the shared memory map. This value is then communicated to the fuzzer and used to allocate exactly the required memory before the fuzzing loop begins.

\paragraph{Coverage Callback} The core of the tracking mechanism is the callback injected into the IR. This function is executed immediately prior to every indirect control flow transfer and, consequently, is optimised for minimal latency.

The indirect map utilises the unique slot ID as a direct index into the map array. However, a single indirect branch source often dispatches execution to multiple distinct target addresses. To record these diverse destinations, the runtime must efficiently map the target address value to a specific bit within the source's Bloom filter.

To achieve this, each destination is associated with a single bit in the source's Bloom filter using a multiplicative hash function applied to the lower 12 bits of the destination address. The hash result is then truncated to yield an index, and the bit at the corresponding position in the Bloom filter is set to 1.

\section{Execution Evaluation}

The introduction of the secondary indirect coverage map necessitates adjustments to the fuzzer's core heuristics. Specifically, it calls for changes in how the engine evaluates the utility of newly discovered test cases and how the execution queue is managed to prioritise inputs that uncover novel indirect branch targets. Substantial modifications to the execution scheduling algorithm are therefore needed.

\paragraph{Queue Management Integration} The cornerstone of CGF is the ability to recognize when an execution trace reaches novel execution paths. In the standard AFL++ architecture, this is solely governed by comparing the primary edge map against a historical record of all seen edges.

With the introduction of the indirect tracking mechanism, this evaluation phase --- responsible for determining whether a test case should be appended to the fuzzer's queue --- must become dual-faceted. A specialized evaluation function is added to the fuzzer's post-execution routine, performing a comparison between the newly generated indirect map and its historical record.

If novel bits are detected in the secondary map, representing a previously unreached indirect target, the input is immediately qualified as "interesting", similarly to the discovery of a new edge in the standard coverage map. This designation of novelty applies even if the primary edge map yielded no new discoveries. This is crucial as it allows the fuzzer to preserve and explore paths that would otherwise have been obscured.

\paragraph{Seed Scoring} Appending a seed to the queue is only the first step. The scheduler must also determine the quality of the input. AFL++'s scheduling algorithm relies on the concept of "favoured" entries, which are selected more frequently for mutation. To minimise corpus size while maintaining maximum coverage, the fuzzer prioritises the fastest and smallest inputs that cover each known edge (as seen in Section \ref{subsec:qandsched}).

To mirror the standard logic for the indirect map, a secondary array of top-rated entries is introduced, specifically for entries valuable in terms of bits set in the secondary map.

The fuzzer evaluates the current test case's characteristics: 
\begin{itemize}
    \item Whether the test case is the very first to trigger a specific bit in the secondary map.

    \item Whether the test case triggers an already known target but represents a more performant execution path than the previously recorded best input for that target.
\end{itemize}
In either case, the input captures the "top-rated" slot for that specific indirect bit. This strict preservation logic ensures the engine explicitly retains the most efficient inputs required to reach every discovered indirect destination, forming a solid baseline for subsequent mutation cycles.

\paragraph{Performance Multipliers} When selecting a test case from the queue to undergo mutation, the fuzzer calculates a "performance score", which determines the intensity and duration of the fuzzing cycle allocated to that seed. Higher scores translate to more mutation cycles. 

The calculation of this score has been augmented to incorporate the density of the indirect coverage map. Testcases that successfully uncover a significant number of paths containing indirect branches are awarded higher performance multipliers. This heuristic guides the mutation engine to focus its computational resources on inputs that have already demonstrated the capability to explore more traditionally obscure areas of the application's CFG. 

\paragraph{Trace Verification and Execution Checksums} To maintain the integrity of the fuzzing state, AFL++ validates execution traces to ensure determinism and identify stability issues. A non-deterministic path can degrade the fuzzer's efficiency and introduce noise into the feedback loop. 

During operations such as test case calibration and trace minimisation, the fuzzer computes checksums of the coverage bitmap to verify stability. The checksumming routines have been modified to incorporate the secondary map.

The fuzzer now computes a combined checksum representing the state of both the primary edge map and the secondary indirect map. If the indirect map diverges between identical runs, the fuzzer's stability metrics will correctly flag the execution as variable. This validation ensures that the fuzzer does not erroneously optimise or prioritise inputs based on unstable indirect control flow, thereby maintaining the fidelity of the overall fuzzing campaign.